\chapter{Methodology and Implementation}
\label{ch:methodology}

This chapter outlines the design and implementation of the proposed methodology
for multi-label \gls{cxr} classification.
A high level overview of the pipeline implemented is presented
in Figure~\ref{fig:design-pipeline}.

\begin{figure}[tb]
    \centering
    \includegraphics[width=\textwidth]{figures/design-pipeline}
    \caption{Design pipeline for the proposed methodology.}
    \label{fig:design-pipeline}
\end{figure}

\section{Dataset and Data Preparation}
\label{sec:dataset-preparation}

As outlined in the experimental design, the CheXpert and NIH ChestX-ray14
datasets were loaded using the \texttt{torchxrayvision} library. This library provides a
standardized interface that handles the complexities of these datasets, such as
mapping labels to a common set of conditions and applying necessary image
transformations.

The CheXpert dataset presented a unique challenge due to its multi-class
labeling system, where each finding could be marked as present, absent, or
uncertain. The handling of these uncertain labels is non-trivial and can
significantly impact model performance. The original CheXpert researchers
systematically evaluated several mapping strategies, including:

\begin{itemize}
    \item \textbf{U-Zeros}: Treating all uncertain labels as negative (0).
    \item \textbf{U-Ones}: Treating all uncertain labels as positive (1).
    \item \textbf{U-Ignore}: Ignoring samples with uncertain labels during training.
\end{itemize}

For this project, the `U-Zeros' approach was implemented, mapping all uncertain
labels to `absent'. This decision provides a conservative and reproducible
baseline, enabling a clear focus on the primary research goal: evaluating the
performance and guarantees of Conformal Risk Control (CRC) and Learn then Test
(LTT). It is acknowledged, however, that this is a simplification, and a comprehensive
future study could investigate how these different label-handling
philosophies interact with the guarantees and outputs produced by these
calibration frameworks.

\section{Model Architectures}

Data from these datasets were split into training, calibration, and test sets,
also using the \texttt{torchxrayvision} library.

The models for this project were chosen to provide a robust and diverse
evaluation of CRC's performance across different modern convolutional neural
network (CNN) architectures. All models were sourced from the torchxrayvision
library to ensure consistency and reproducibility, and care was taken to ensure
that any results were in line with the library's benchmarks.

ResNet50: A 50-layer Residual Network was selected as it is a foundational
architecture in modern deep learning. Its use of residual connections allows for
the effective training of very deep networks by mitigating the vanishing
gradient problem. It serves as a strong and widely used baseline for
classification tasks.

DenseNet121: A 121-layer Densely Connected Convolutional Network was chosen to
represent an alternative architectural philosophy. DenseNet connects each layer
to every other layer in a feed-forward fashion, promoting feature reuse,
improving feature propagation, and often achieving high performance with fewer
parameters than comparable ResNet models.

Official CheXpert Model: To ensure a direct and fair comparison on the CheXpert
dataset, the official baseline model provided by the dataset's authors (a
DenseNet121 variant) was also included. This model is specifically pre-trained
and configured for the CheXpert task.

The remaining portions of the pipeline were implemented from scratch using
standard PyTorch and NumPy libraries.
